\section{离线算法}

\subsection{莫队算法 (Mo's Algorithm)}

\subsubsection{普通莫队 (Basic Mo's Algorithm)}

\textbf{用途：}
离线处理大量\textbf{静态区间询问}：若所求统计量能以 $O(1)$ 代价在相邻区间 $[l\pm1,r]$、$[l,r\pm1]$ 之间增删推得（即单个元素进出区间时答案可增量维护），则通过合理排序询问，用双指针在数组上滑动，总移动量从暴力的 $O(nm)$ 降到 $O(n\sqrt{m})$ 级别。示例问题：$m$ 次询问区间 $[l,r]$ 内\textbf{不同数的个数}。

\textbf{核心概念：}
\begin{itemize}
  \item \textbf{离线：} 询问一次性读入、排序后统一回答，最后按原序输出；不支持强制在线。
  \item \textbf{分块排序：} 下标 $1..n$ 按块长 $B$ 分块，询问按（左端点所在块号，右端点）双关键字排序。
  \item \textbf{双指针增量维护：} 维护当前区间与答案，指针每移动一步，仅对新进/离开区间的元素执行 \texttt{add}/\texttt{del}。
  \item \textbf{奇偶优化：} 奇数块内右端点升序、偶数块内降序，右端点在块间切换时不必先跳回左端，实测约省一半移动量。
\end{itemize}

\textbf{算法流程：}
\begin{enumerate}
  \item 读入数组并离散化，读入全部询问，取块长 $B = n/\sqrt{m}$。
  \item 询问排序后，从空区间 $l=1, r=0$ 出发，依次将指针移至每个询问区间，先扩张后收缩，边移边 \texttt{add}/\texttt{del}。
  \item 指针就位后记录当前答案，全部处理完按原序输出。
\end{enumerate}

\textbf{时间复杂度：} 左端点同块内每步至多移动 $B$，共 $O(mB)$；右端点在每块内单调不降，共 $n/B$ 块，合计 $O(n^2/B)$。总复杂度 $O(mB + n^2/B)$，取 $B=n/\sqrt{m}$ 时为 $O(n\sqrt{m})$。

\lstinputlisting[language=C++]{codes/08-offline-01.cpp}

\begin{quote}
\textbf{注意：} 换题只需修改 \texttt{add}/\texttt{del}（元素下标进出区间）与答案变量的维护方式，排序与指针骨架完全不变。统计量必须\textbf{支持增量删除}——区间最值这类无法 $O(1)$ 撤销的量需改用回滚莫队、树上询问需树上莫队等变体。
\end{quote}

\subsubsection{带修莫队 (Mo's Algorithm with Updates)}

\textbf{用途：}
在普通莫队基础上支持\textbf{单点修改}与区间询问穿插：把「时间」（已执行的修改个数）作为第三维关键字，与左右端点一起排序、一起移动。示例问题：单点改值 + 询问区间内不同数的个数。

\textbf{核心概念：}
\begin{itemize}
  \item \textbf{三维排序：} 询问按（左端点所在块，右端点所在块，时间戳）三关键字升序排序。
  \item \textbf{时间指针：} 维护当前已执行的修改数 $T$，处理每个询问前把修改序列推进或回退到询问对应的时刻。
  \item \textbf{swap 撤销：} 执行第 $i$ 个修改时 \texttt{swap(a[p], upd[i].val)}，\texttt{val} 由新值变为旧值；再 swap 一次即还原，故执行与撤销共用同一段代码。
  \item 修改位置落在当前区间内时，先 \texttt{del} 旧值再 \texttt{add} 新值（撤销时同理，方向对称）。
\end{itemize}

\textbf{算法流程：}
\begin{enumerate}
  \item 离线读入全部操作，修改与询问分别存储；每个询问记录其之前的修改个数 $t$。初始值与全部修改值一起离散化。
  \item 取块长 $B = n^{2/3}$，按三关键字排序。
  \item 依次移动 $l, r, T$ 三个指针：时间指针前移为执行修改、后移为撤销修改，就位后记录答案。
\end{enumerate}

\textbf{时间复杂度：} 三指针总移动量在询问数、修改数与 $n$ 同阶时为 $O(n^{5/3})$（块长 $n^{2/3}$）。

\lstinputlisting[language=C++]{codes/08-offline-02.cpp}

\begin{quote}
\textbf{注意：} 块长对常数影响显著，卡常时可微调 $B$。离散化必须包含\textbf{所有修改值}，否则后续修改会引入未映射的新值。
\end{quote}
